% Automatically rendered from a complete C07 summary; no new statistics.
% Summary SHA256: 909ef28817342c19293fe9cb832bcf659dd480be5fd82c3708ee73b832cd4f9a
\begin{table}[htbp]
\centering
\caption{Resolução numérica dos PITs: número de funções resolvidas entre os 500 dados de cada modelo.}
\label{tab:c07-resolucao-producao}
\begingroup\small\setlength{\tabcolsep}{4pt}
\begin{tabular}{@{}lrrrrrrr@{}}
\toprule
Modelo & $u$ & $\log A$ & $\gamma$ & $\log A_r$ & $\log E$ & $\log L$ & Todas \\
\midrule
A0\_CN & 497 & 497 & 497 & 497 & 497 & 497 & 497 \\
A\_CN & 490 & 487 & 489 & 489 & 489 & 488 & 485 \\
B\_CN & 495 & 495 & 495 & 495 & 495 & 494 & 494 \\
A\_G & 489 & 487 & 489 & 489 & 487 & 488 & 484 \\
B\_G & 490 & 491 & 491 & 491 & 491 & 491 & 490 \\
\bottomrule
\end{tabular}
\endgroup
\par\smallskip\begin{minipage}{0.96\textwidth}\footnotesize As colunas seguem os cinco parâmetros e a log-verossimilhança na verdade. Cada função é avaliada pelos seus próprios diagnósticos; funções não resolvidas conservam o intervalo $[0,1]$. A última coluna exige os seis PITs resolvidos na mesma realização e não aprova todos os vinte quantis.\end{minipage}
\end{table}
\begin{table}[htbp]
\centering
\caption{Contagem de rejeições entre os 31 testes de cada modelo, com ajuste de Holm a 5\%.}
\label{tab:c07-testes-producao}
\begingroup\small\setlength{\tabcolsep}{4pt}
\begin{tabular}{@{}lrrrr@{}}
\toprule
Modelo & Família & Nominais & Robustas & Possíveis \\
\midrule
A0\_CN & 93 & 0 & 0 & 0 \\
A\_CN & 62 & 8 & 6 & 18 \\
B\_CN & 62 & 5 & 4 & 5 \\
A\_G & 93 & 0 & 0 & 1 \\
B\_G & 93 & 0 & 0 & 0 \\
\bottomrule
\end{tabular}
\endgroup
\par\smallskip\begin{minipage}{0.96\textwidth}\footnotesize As famílias de 93 e 62 hipóteses são separadas. Uma rejeição robusta persiste para todos os PITs dos envelopes numéricos; uma rejeição possível não é excluída por esses envelopes. As contagens são testes dependentes, não descobertas independentes. Não rejeitar não demonstra correção.\end{minipage}
\end{table}
\begin{table}[htbp]
\centering
\caption{Teste de Kolmogorov--Smirnov para o PIT da massa adimensional $u$.}
\label{tab:c07-pit-massa}
\begingroup\small\setlength{\tabcolsep}{4pt}
\begin{tabular}{@{}lrrrr@{}}
\toprule
Modelo & $D_{\rm KS}$ & $p_{\rm Holm}$ & Faixa de $p_{\rm Holm}$ & Resolvidos \\
\midrule
A0\_CN & $0{,}038$ & $1$ & $[1;1]$ & 497 \\
A\_CN & $0{,}033$ & $1$ & $[1;1]$ & 490 \\
B\_CN & $0{,}031$ & $1$ & $[1;1]$ & 495 \\
A\_G & $0{,}035$ & $1$ & $[1;1]$ & 489 \\
B\_G & $0{,}033$ & $1$ & $[1;1]$ & 490 \\
\bottomrule
\end{tabular}
\endgroup
\par\smallskip\begin{minipage}{0.96\textwidth}\footnotesize O teste nominal condiciona-se aos PITs calculados. A faixa é uma análise de sensibilidade à margem $0{,}002+z\,\mathrm{MCSE}$ e aos intervalos $[0,1]$ das funções não resolvidas; não é um intervalo de confiança exato para o valor-$p$. A margem determinística é operacional, não uma garantia uniforme da integral.\end{minipage}
\end{table}
\begin{table}[htbp]
\centering
\caption{Coberturas da massa sob a priori e sua sensibilidade numérica.}
\label{tab:c07-cobertura-massa}
\begingroup\small\setlength{\tabcolsep}{4pt}
\begin{tabular}{@{}llrrr@{}}
\toprule
Modelo & Evento & Contagem & IC binomial 95\% & Fração possível \\
\midrule
A0\_CN & $\mathrm{PIT}_u\leq0{,}95$ & 482/500 & $[0{,}944;0{,}979]$ & $[0{,}956;0{,}964]$ \\
A\_CN & $\mathrm{PIT}_u\leq0{,}95$ & 480/500 & $[0{,}939;0{,}975]$ & $[0{,}936;0{,}962]$ \\
B\_CN & $\mathrm{PIT}_u\leq0{,}95$ & 478/500 & $[0{,}934;0{,}972]$ & $[0{,}946;0{,}966]$ \\
A\_G & $\mathrm{PIT}_u\leq0{,}95$ & 478/500 & $[0{,}934;0{,}972]$ & $[0{,}930;0{,}962]$ \\
B\_G & $\mathrm{PIT}_u\leq0{,}95$ & 480/500 & $[0{,}939;0{,}975]$ & $[0{,}934;0{,}964]$ \\
A0\_CN & $0{,}05\leq\mathrm{PIT}_u\leq0{,}95$ & 460/500 & $[0{,}893;0{,}942]$ & $[0{,}912;0{,}924]$ \\
A\_CN & $0{,}05\leq\mathrm{PIT}_u\leq0{,}95$ & 456/500 & $[0{,}884;0{,}935]$ & $[0{,}886;0{,}916]$ \\
B\_CN & $0{,}05\leq\mathrm{PIT}_u\leq0{,}95$ & 455/500 & $[0{,}881;0{,}934]$ & $[0{,}898;0{,}922]$ \\
A\_G & $0{,}05\leq\mathrm{PIT}_u\leq0{,}95$ & 457/500 & $[0{,}886;0{,}937]$ & $[0{,}888;0{,}926]$ \\
B\_G & $0{,}05\leq\mathrm{PIT}_u\leq0{,}95$ & 458/500 & $[0{,}888;0{,}939]$ & $[0{,}890;0{,}922]$ \\
\bottomrule
\end{tabular}
\endgroup
\par\smallskip\begin{minipage}{0.96\textwidth}\footnotesize O intervalo binomial de Clopper--Pearson é ponto a ponto e condiciona-se à contagem nominal; não é simultâneo entre modelos ou eventos. A última coluna representa a faixa de frações compatíveis com os PITs numéricos. As uniões de intervalos binomiais correspondentes estão no JSON completo. Eventos são avaliados diretamente pela CDF, sem supor erro horizontal nulo nos quantis.\end{minipage}
\end{table}
